
\begin{abstract}

Simultaneous speech translation produces target text incrementally from partial speech input.
Recent speech large language models have markedly improved SST quality but still struggle with rare and domain-specific terminology. Retrieval augmentation has helped in automatic speech recognition and neural machine translation, but extending it to SST is non-trivial: retrieval must be fast and accurate under partial speech, and the model must decide whether and when to apply retrieved terms during incremental generation.
We propose \textbf{R}etrieval-\textbf{A}ugmented \textbf{S}imultaneous \textbf{S}peech \textbf{T}ranslation (\method), which addresses both challenges. 
For accurate cross-modal retrieval under partial input, \method trains a lightweight speech–text retriever that produces chunkwise terminology hints for the Speech LLM via multi-scale retrieval. To use these hints correctly, we synthesize training data that teaches the Speech LLM to decide whether and when to apply each retrieved term. \crchange{On the human-annotated ACL 60/60 dev set, \method improves exact-form terminology accuracy over InfiniSST in every tested language and latency setting, by 12.3--21.4 percentage points, while BLEU and XCOMET show mixed but mostly positive changes. On the human-reference En--De portion of the ESO test set, the terminology gains are 26.0--31.2 points. The retriever adds negligible computational overhead.}

\end{abstract}
